\documentclass[11pt,a4paper]{article}
\usepackage[utf8]{inputenc}
\usepackage[spanish]{babel}
\usepackage{amsmath}
\usepackage{amsfonts}
\usepackage{amssymb}
\usepackage{geometry}
\usepackage{enumerate}
\geometry{top=2.5cm, bottom=2.5cm, left=2.5cm, right=2.5cm}

\begin{document}

\begin{center}
    \textbf{\Large SEGUNDO PARCIAL DE FÍSICA BÁSICA} \\[0.2cm]
    \textbf{\large VERSIÓN B} \\[0.4cm]
\end{center}

\noindent \textbf{Estudiante:} \hrulefill \ \textbf{Matrícula:} \underline{\hspace{3cm}} \\[0.2cm]
\textbf{Fecha:} \underline{\hspace{4cm}} \ \textbf{Calificación:} \underline{\hspace{1.5cm}} / 8.0 puntos \\

\hrule
\section*{I. Teoría (Valor: 4.0 puntos - 0.5 puntos cada una)}
\textit{Selecciona la opción correcta en cada una de las siguientes preguntas. No debe haber tachaduras ni enmiendas.}

\begin{enumerate}
    \item \textbf{De acuerdo con la Primera Ley de Newton, si la fuerza neta que actúa sobre un cuerpo es nula, el cuerpo necesariamente:}
    \begin{enumerate}[a)]
        \item Debe estar en reposo absoluto.
        \item Permanecerá en reposo o en movimiento rectilíneo uniforme.
        \item Experimentará un aumento constante de su velocidad.
        \item Cambiará continuamente su dirección de movimiento.
    \end{enumerate}

    \item \textbf{De acuerdo con la Tercera Ley de Newton (Acción-Reacción), si un cuerpo A ejerce una fuerza $\vec{F}_{AB}$ sobre un cuerpo B, entonces:}
    \begin{enumerate}[a)]
        \item El cuerpo B ejerce sobre A una fuerza de igual magnitud, misma dirección y sentido contrario ($\vec{F}_{BA} = -\vec{F}_{AB}$).
        \item El cuerpo B ejerce sobre A una fuerza de mayor magnitud en el mismo sentido.
        \item La fuerza del cuerpo B sobre A es nula porque A inició el contacto.
        \item El cuerpo B ejerce sobre A una fuerza perpendicular a la original.
    \end{enumerate}

    \item \textbf{Despreciando la resistencia del aire, el ángulo de disparo respecto a la horizontal que produce el máximo alcance horizontal para un proyectil es:}
    \begin{enumerate}[a)]
        \item $30^\circ$
        \item $45^\circ$
        \item $60^\circ$
        \item $90^\circ$
    \end{enumerate}

    \item \textbf{La propiedad física de un cuerpo que mide cuantitativamente su inercia (resistencia al cambio en su estado de movimiento) es la:}
    \begin{enumerate}[a)]
        \item Aceleración.
        \item Fuerza normal.
        \item Masa.
        \item Fricción.
    \end{enumerate}

    \item \textbf{Con respecto a la fuerza de fricción entre dos superficies secas en contacto, se cumple que:}
    \begin{enumerate}[a)]
        \item El coeficiente de fricción cinética ($\mu_c$) es siempre menor que el coeficiente de fricción estática ($\mu_e$).
        \item La fuerza de fricción estática tiene un valor constante e invariable.
        \item La fuerza de fricción es siempre perpendicular a la superficie de contacto.
        \item El coeficiente de fricción depende del área de contacto visible.
    \end{enumerate}

    \item \textbf{La fuerza es una magnitud física de carácter:}
    \begin{enumerate}[a)]
        \item Escalar.
        \item Vectorial.
        \item Adimensional.
        \item Tensorial de segundo orden.
    \end{enumerate}

    \item \textbf{A diferencia de la masa, el peso de un objeto:}
    \begin{enumerate}[a)]
        \item Es una propiedad intrínseca que no varía con la gravedad local.
        \item Es la fuerza de atracción gravitatoria que ejerce la Tierra sobre él ($w = m \cdot g$).
        \item Se mide en kilogramos en el Sistema Internacional.
        \item Es una cantidad escalar pura sin dirección ni sentido.
    \end{enumerate}

    \item \textbf{La Segunda Ley de Newton establece matemáticamente que la aceleración de un cuerpo es:}
    \begin{enumerate}[a)]
        \item Directamente proporcional a su masa e inversamente proporcional a la fuerza neta.
        \item Directamente proporcional a la fuerza neta e inversamente proporcional a su masa.
        \item Independiente de la masa y proporcional al peso.
        \item Proporcional al cuadrado de la fuerza neta aplicada.
    \end{enumerate}

\end{enumerate}

\hrule
\section*{II. Problemas Prácticos (Valor: 4.0 puntos - 1.0 punto cada uno)}
\textit{Resuelve cada uno de los problemas y selecciona la respuesta correcta. Muestra tus cálculos breves.}

\begin{enumerate}
    \item \textbf{Un cuerpo pesa $14.7 \text{ N}$ en la Tierra ($g = 9.80 \text{ m/s}^2$). ¿Cuánto pesará este mismo cuerpo en la superficie de la Luna, donde la aceleración de la gravedad es de $1.62 \text{ m/s}^2$?:}
    \begin{enumerate}[a)]
        \item $14.7 \text{ N}$
        \item $2.43 \text{ N}$
        \item $23.8 \text{ N}$
        \item $1.50 \text{ N}$
    \end{enumerate}

    \item \textbf{Un cuerpo de masa $m = 3.00 \text{ kg}$ acelera horizontalmente a razón de $a_x = +0.800 \text{ m/s}^2$ (hacia la derecha). Si actúan dos fuerzas sobre el eje $x$ y una de ellas es $F_{2x} = +9.00 \text{ N}$, determine la componente de la otra fuerza $F_{1x}$:}
    \begin{enumerate}[a)]
        \item $+6.60 \text{ N}$
        \item $-6.60 \text{ N}$
        \item $-11.4 \text{ N}$
        \item $-2.40 \text{ N}$
    \end{enumerate}

    \item \textbf{Una bola de golf se desliza por el césped con una velocidad inicial de $6.00 \text{ m/s}$ y llega al hoyo a $2.00 \text{ m/s}$ después de recorrer una distancia de $10.0 \text{ m}$ en línea recta. Determine su aceleración constante:}
    \begin{enumerate}[a)]
        \item $-1.60 \text{ m/s}^2$
        \item $+1.60 \text{ m/s}^2$
        \item $-3.20 \text{ m/s}^2$
        \item $-0.80 \text{ m/s}^2$
    \end{enumerate}

    \item \textbf{Un estudiante deja caer un tomate (desde el reposo) desde la azotea de un edificio y mide que tarda exactamente $1.60 \text{ s}$ en golpear el suelo. Despreciando la resistencia del aire y usando $g = 9.80 \text{ m/s}^2$, calcule la altura del edificio:}
    \begin{enumerate}[a)]
        \item $15.7 \text{ m}$
        \item $12.5 \text{ m}$
        \item $25.1 \text{ m}$
        \item $7.84 \text{ m}$
    \end{enumerate}

\end{enumerate}

\end{document}
